\documentclass{beamer}
\usepackage{verbatim}
\usepackage{wrapfig}
\usepackage{color} 
\usetheme{Warsaw}
\title{A gentle intro to Golang and the Go-universe}
\author{Alex-P. Natsios}
\institute{GoCode.Thessaly(7)}
\date{28 feb 2014}
\begin{document}
    \begin{frame}
       \titlepage
    \end{frame}

    \begin{frame}
        \center\huge The Language
    \end{frame}

    \begin{frame}{A few words about GO}

        {\bf Go}, also called {\bf golang}, was initially developed at google in 2007 and
        announced in 2009 but most of its popularity boost (and hype) came
        much later with its first stable release (1.0) on 28 Mar 2012.

        \begin{block}{Characteristics:}
            \begin{itemize}
                \item Compiled
                \item Statically typed and garbage collected
                \item Object Oriented (but not in the usual way)
                \item Sane Concurrency
                \item Fast Compilers
                \item Rich Standard Library
                \item Scalable Tools
            \end{itemize}
        \end{block}
    \end{frame}

    \begin{frame}{The Gopher!}
        \begin{figure}{l}{40mm}
            \includegraphics[scale=0.5]{images/gopher.eps}
        \end{figure}
    \end{frame}

    \begin{frame}{The Classic Example}
        \verbatiminput{code/hello.go}
    \end{frame}

    \begin{frame}{Packages}
        \begin{itemize}
            \item Go code lives in packages.\\
            \item A package consists of one or more source(.go) files.\\
            \item Typically all files belonging to a package are placed in the same directory.
            \item Visibility is determined by case: Foo is exported, foo is not.
        \end{itemize}
    \end{frame}

    \begin{frame}{Declarations}
        import ``fmt''\\
        const answer = 42\\
        var something [10]byte\\
        var names = []string \textbraceleft ``Alex'', ``George'', ``Maria'' \textbraceright \\
    \end{frame}

    \begin{frame}{Types}
        \begin{itemize}
            \item All basic types you would expect to have.\\
            \item But string is a basic type.\\
            \item No pointer arithmetic.\\
            New (or not exactly that new) types:\\
            \item Slices: Much like many other modern languages, []int \\
            \item Maps: Because suddenly everybody needs them, map[string]int \\
            \item Interfaces: for Polymorphism, interface\textbraceleft \textbraceright \\
            \item Channels: Used to communicate with goroutines, chan int
        \end{itemize}
    \end{frame}

    \begin{frame}{Scope}
        \begin{itemize}
            \item Universe
            \item Package
            \item File (for imports)
            \item Function
            \item Block
        \end{itemize}
    \end{frame}

    \begin{frame}{Concurrency}
       Concurrency is a property of systems in which several computations are
       executing simultaneously and potentialy interacting with each other.
       
       \emph{ {\color{red}WARNING:}} Concurrency is {\color{red}NOT} parallelism,
       although it enables parallelism if you have a multiprocessor system.
    \end{frame}

    \begin{frame}{Goroutines}
        \begin{itemize}
            \item Independently executing functions, launched by a {\bf go} statement. 
            \item \emph{ {\color{red}NOT}} threads.
            \item VERY cheap, you might have thousands of goroutines running
                under the same thread.
            \item Have their own dynamic call stack (growning and shrinking as needed).
        \end{itemize}
    \end{frame}

    \begin{frame}{Channels}
        A channel provides a connection between goroutines, allowing communication.\\
        Channels can be unbuffered or buffered, so they both communicate and synchronize.\\
        Buffered channels are asynchronous.
    \end{frame}

    \begin{frame}{Channels}{Code Example}
        \verbatiminput{code/channels.go}
    \end{frame}

    \begin{frame}{Channels}{cont.}
        When in main(), ``\textless-greetings'' is executed it waits for a value to be sent.\\
        Same goes for our anonymous function that expects a reciever to be in place
        in order for the greetings to be sent.\\
        If no sender/receiver is ready (they both must be) then we wait until they are!
    \end{frame}

    \begin{frame}
        \center\huge The Toolchain
    \end{frame}

    \begin{frame}{The go tool}

        {\bf go build} - To compile the package.\\
        {\bf go get} - To resolve and install deps.\\
        {\bf go test} - To run the test suite and benchmarks.\\
        {\bf go install} - To install the package.\\
        {\bf go doc} - To generate documentation.\\
        {\bf go fmt} - To properly format your code.\\
        {\bf go run} - To build and run the app.\\
        {\bf go tool} - To run extra tools.\\
        and more (properly ``integrated'' in the {\bf go} tool)

    \end{frame}

    \begin{frame}{Building and workspaces}
        A {\bf GO} program can be compiled and linked without additional build info.
        A single tool can compile either individual files or entire systems.

        In order to work without build scripts a certain directory structure
        {\bf MUST} be followed.
        \begin{block}{workspace}
            workspace/\\
            \ \ bin/\\
            \ \ pkg/\\
            \ \ src/\\
        \end{block}
    \end{frame}

    \begin{frame}{workspaces}{cont.}
        \begin{block}{Creating a workspace:}
            mkdir -p \textdollar HOME/GoCode/\textbraceleft bin, pkg, src\textbraceright
        \end{block}

        \begin{block}{Telling go about it:}
            export GOPATH=``\textdollar HOME/GoCode''\\
            export PATH=``\textdollar PATH:\textdollar GOPATH/bin''
        \end{block}
    \end{frame}

    \begin{frame}{Links}
        Golang homepage:\\
        \ \ \url{golang.org}\\
        Go tour:\\
        \ \ \url{tour.golang.org}\\
        Package Doc:\\
        \ \ \url{golang.org/pkg}\\
        A little outdated but still useful page:\\
        \ \ \url{go-lang.cat-v.org}\\
        Obligatory subreddit:\\
        \ \ \url{reddit.com/r/golang/}
    \end{frame}

    \begin{frame}{Thank you for your attention!}
        \center Alex-P. Natsios\\
        \center\url{drakevr@2f30.org}\\
        \center\url{http://drakevr.gr}\\
        \center\url{http://www.linkedin.com/in/drakevr}\\
        \center\url{http://www.github.com/drakevr}\\
        \center\url{http://www.facebook.com/drakevr}\\
        \center\url{http://www.twitter.com/drakevr}
    \end{frame}

\end{document}
